% Part: first-order-logic
% Chapter: introduction
% Section: first-order-logic

\documentclass[../../../include/open-logic-section]{subfiles}

\begin{document}

\olfileid{fol}{int}{fol}

\olsection{د لومړۍ درجې منطق}

تاسو غالباً د صوري منطق په لومړۍ پېژندنه کښې د لومړۍ درجې له منطق
سره بلد يئ۔\footnote{په حقيقت کښې موږ لږ و ډېر دا فرضوو چې تاسو
ورسره بلد يئ!{} که نۀ يئ، نو يو لا ابتدايي درسي کتاب، لکه
\emph{\texttt{forall x}} \citep{Magnus2021}، بيا کتلے شئ۔} ښايي تاسو دا د
``کميت ټاکونکو منطق'' يا ``پريديکاتي منطق'' په نامه پېژنئ۔ د لومړۍ
درجې منطق تر هر څه مخکښې يوه صوري ژبه ده۔ يعنې يو ټاکلے لغتځونډ
لري، او عبارتونه ئې د همدې لغتځونډ تارونه دي۔ خو هر تار جايز نۀ
دے۔ د جايزو عبارتونو بېلابېل ډولونه شته: ترمونه، !!{formula}s او
!!{sentence}s۔ موږ په اصل کښې د لومړۍ درجې منطق له !!{sentence}s
سره لېوال يو: هغوئ د انګرېزي ژبې د جملو صوري سيال راکوي، او د هغوئ
په باب هغه پوښتنې کولے شو چې منطقپوه عموماً ورسره لېوال وي۔ مثلاً:
\begin{itemize}
    \item ايا $!B$ له~$!A$ څخه په منطقي ډول لازمېږي؟
    \item ايا $!A$ په منطقي ډول رښتونې، په منطقي ډول دروغجنه، که
اتفاقي ده؟
    \item ايا $!A$ او $!B$ هم‌ارزښته دي؟
\end{itemize}

دا پوښتنې په بنسټيز ډول د لومړۍ درجې منطق د !!{sentence}s د
``معنٰی'' په باب دي۔ مثلاً يو فيلسوف به دا پوښتنه، چې ايا $!B$ له
~$!A$ څخه په منطقي ډول لازمېږي، داسې وڅېړي: ايا کوم داسې حالت شته
چې $!A$ پکښې رښتونې خو~$!B$ دروغجنه وي ($!B$ له~$!A$ څخه نۀ
لازمېږي)، که هر هغه حالت چې $!A$ رښتونې کوي~$!B$ هم رښتونې کوي
($!B$ له~$!A$ څخه لازمېږي)؟ خو لا راته نۀ دي ويل شوي چې ``حالت''
څۀ دے---دا د \emph{معنٰی پوهنې} کار دے۔ د لومړۍ درجې منطق معنٰی
پوهنه د فيلسوف د ``حالت'' د وجداني مفکورې، او هم---چې دا مهمه
ده---په يوه حالت کښې د !!a{sentence}~$!A$ د \emph{رښتونې کېدو}
رياضيکي دقيق مدل راکوي۔ هغه رياضيکي دقيق مدل چې جوړوو ئې
!!a{structure} بولو۔ هغه اړيکه چې ``پکښې رښتونې'' په دقيق ډول
بيانوي، د \emph{صدق} اړيکه بلل کېږي۔ نو موږ به د
!!{sentence}s~$!A$ او !!{structure}s~$\Struct{M}$ دپاره دا تعريفوو
چې ``$!A$ په~$\Struct{M}$ کښې صدق کوي'' (په سمبولونو کښې:
$\Sat{M}{!A}$)۔ له دې وروسته نور معنايي اصطلاحات، لکه ``ترې
لازمېږي'' يا ``په منطقي ډول رښتونې ده''، هم په دقيق ډول تعريفولے
شو۔ دا تعريفونه به ممکنه کړي چې په رياضيکي دقت سره، مثلاً، فيصله
وکړو چې ايا
$\lforall[x][(!A(x) \lif !B(x))], \lexists[x][!A(x)] \Entails \lexists[x][!B(x)]$
ده که نۀ۔ ځواب به طبعاً ``هو'' وي۔
\textbf{د ژباړې د نسخې سمون (OLFOL-001):} په سرچينه کښې د دې
استلزام د لومړۍ عمومي جملې تړونکى قوس تر وروستۍ پايلې پورې ځنډېدلے
او کميت ټاکونکى ئې په ناسم ډول پر ټول استلزام خور کړے و؛ دلته قوس
د شرطي فارمول سمدستي وروسته تړل شو۔ که تاسو مخکښې د
انګرېزي جملو د لومړۍ درجې منطق په سمبولونو اړولو زده کړه کړې وي، دا
به مثلاً د ``ټولې مېږي حشرات دي، مېږي شته، نو حشرات شته'' د
سمبولونو بڼه وپېژنئ۔ دا په څرګند ډول معتبر استدلال دے؛ نو زموږ د
صوري ژبې د ``ترې لازمېږي'' رياضيکي مدل هم بايد همدا ځواب ورکړي۔

بله موضوع چې غالباً د صوري منطق له لومړۍ پېژندنې څخه مو ياده ده،
\emph{!!{derivation}s} دي۔ که مو د صوري منطق لومړے کورس لوستلے وي،
ښوونکي به درباندې د داسې !!{derivation}s د موندلو تمرين کړے وي؛
ښايي حتٰی داسې !!a{derivation} مو موندلے وي چې پورته استلزام
ښيي۔ د !!{derivation}s د ورکولو ډېرې بېلابېلې لارې شته: ښايي تاسو
``فطري استنتاج'' يا ``د رښتياوالي ونې'' کارولې وي، خو نورې ډېرې
لارې هم شته۔ د !!{derivation} د نظامونو موخه داسې وسيلې برابرول دي
چې د منطقپوهانو پورته پوښتنې پرې ځواب شي: مثلاً د فطري استنتاج داسې
!!{derivation} چې $\lforall[x][(!A(x) \lif !B(x))]$ او
$\lexists[x][!A(x)]$ ئې مقدمې وي او $\lexists[x][!B(x)]$ ئې پايله
(وروستۍ کرښه) وي، دا \emph{تصديقوي} چې $\lexists[x][!B(x)]$ په
منطقي ډول له $\lforall[x][(!A(x) \lif !B(x))]$ او
$\lexists[x][!A(x)]$ څخه لازمېږي۔ په دې پاراګراف کښې هم د
OLFOL-001 له مخې د هماغې عمومي جملې تړونکى قوس او د پايلې زيات قوس
سم شوي دي۔

خو دا ولې؟ په ظاهره د !!{derivation} نظامونه له معنٰی پوهنې سره
هېڅ تړاو نۀ لري: د يوه صوري !!{derivation} ورکول يوازې دا غواړي چې
سمبولونه د ځينو قاعدو له مخې سره منظم شي؛ هلته د ``حالتونو'' يا
``پکښې رښتونې'' هېڅ يادونه نشته۔ د !!{derivation} د نظامونو او
معنٰی پوهنې تر منځ تړاو بايد په مابعدمنطقي څېړنه ثابت شي۔ پکار دا
ده چې مثلاً په رياضيکي ډول ثابته کړو: له مقدمو
$\lforall[x][(!A(x) \lif !B(x))]$ او $\lexists[x][!A(x)]$ څخه د
$\lexists[x][!B(x)]$ صوري !!{derivation} هله او يوازې هله ممکن دے
چې $\lforall[x][(!A(x) \lif !B(x))]$ او
$\lexists[x][!A(x)]$ په ګډه $\lexists[x][!B(x)]$ لازم کړي۔ د
OLFOL-001 له مخې دلته هم د تکرار شوې شېما دوه ناسمه قوسونه سم شوي
دي۔ خو تر دې مخکښې چې دا کار وشي، د نحو او معنٰی پوهنې د تعريفونو
د سمولو دپاره ډېر زيارکښ کار کول پکار دي۔

\end{document}
